% BÀI TOÁN HÌNH HỌC - BAI_HINH_47 (TRÍCH ĐỀ 47)
\begin{ex}
% Câu 13
	Cho tam giác ABC có ba góc nhọn nội tiếp trong đường tròn $(O;R)$. Các đường cao AD, BE và CF cắt nhau tại H. Gọi M và N thứ tự là giao điểm thứ hai của đường tròn $(O;R)$ với BE và CF.
	\begin{enumerate}[a)]
		\item Chứng minh tứ giác BCEF nội tiếp đường tròn.
		\item Chứng minh $\widehat{NMB} = \widehat{BEF}$; AO cắt EF tại K. Gọi P, Q lần lượt là trung điểm của EF và DK. Chứng minh 3 điểm A, P, Q thẳng hàng.
	\end{enumerate}
\loigiai{
{
}
\end{ex}
